% architecture-infographic.tex — Agent Framework Suite architecture infographic
% Compile: python3 /Users/junix/.codex/skills/create-tikz-infographic/scripts/render_tikz.py \
%   docs/architecture-infographic.tex --out-dir /tmp/tikz-agent-suite --engine xelatex --dpi 300 --svg
\documentclass[border=12pt,varwidth=16cm]{standalone}
\usepackage[UTF8,fontset=none]{ctex}
\IfFontExistsTF{Source Han Serif SC}{
  \setmainfont[BoldFont={Source Han Serif SC Bold}]{Source Han Serif SC}
  \setsansfont[BoldFont={Source Han Serif SC Bold}]{Source Han Serif SC}
  \setCJKmainfont[BoldFont={Source Han Serif SC Bold}]{Source Han Serif SC}
  \setCJKsansfont[BoldFont={Source Han Serif SC Bold}]{Source Han Serif SC}
}{
  \PackageWarningNoLine{tikz-infographic}{Source Han Serif SC not found; using Fandol}
  \setCJKmainfont[BoldFont={FandolSong-Bold}]{FandolSong-Regular}
  \setCJKsansfont[BoldFont={FandolHei-Bold}]{FandolHei-Regular}
}
\IfFontExistsTF{Sarasa Mono SC}{
  \setmonofont{Sarasa Mono SC}
  \setCJKmonofont{Sarasa Mono SC}
}{
  \PackageWarningNoLine{tikz-infographic}{Sarasa Mono SC not found; using TeX defaults}
  \setmonofont{Latin Modern Mono}
  \setCJKmonofont{FandolFang-Regular}
}

\usepackage{xcolor}
\usepackage{tikz}
\usetikzlibrary{arrows.meta,backgrounds,calc,fit,positioning}

% ── Use showcase.sty color and style tokens ─────────────────────────────
\input{styles/showcase}

% ── Project-specific styles ─────────────────────────────────────────────
\tikzset{
  suite/.style={
    showcase/node, fill=Ocean!12, draw=Ocean, line width=1pt,
    minimum width=42mm, minimum height=12mm
  },
  driver/.style={
    showcase/node, fill=Sky!18, draw=Sky, line width=.9pt,
    minimum width=34mm, minimum height=11mm
  },
  framework/.style={
    showcase/node, fill=Fog!60, draw=Slate!55, line width=.8pt,
    minimum width=34mm, minimum height=11mm
  },
  actor/.style={
    showcase/node, fill=PaperDeep, draw=Slate!50, line width=.8pt,
    minimum width=30mm, minimum height=10mm, rounded corners=4mm
  },
  component/.style={
    showcase/node, fill=white, draw=Ocean!40, line width=.8pt,
    minimum width=28mm, minimum height=9mm
  },
  boundary/.style={
    draw=Slate!40, densely dashed, line width=.8pt, rounded corners=3mm, inner sep=5mm
  },
  pill/.style={
    showcase/pill, fill=Mint, text=Ocean
  },
  tag/.style={
    showcase/pill, fill=Sky!20, text=Ocean, minimum height=5mm, inner xsep=2mm
  },
  section label/.style={
    showcase/eyebrow, anchor=west
  },
  section title/.style={
    showcase/title, anchor=west
  },
  section desc/.style={
    showcase/subtitle, anchor=west
  },
  legend item/.style={
    font=\scriptsize, text=Slate
  },
  phase box/.style={
    showcase/card, inner sep=4mm, minimum width=14cm
  }
}

\newcommand{\SectionHeader}[3]{%
  \node[section label] at (0,#1) {#2};
  \pgfmathsetmacro{\tmpy}{#1-0.52}
  \node[section title] at (0,\tmpy) {#3};
}

\begin{document}
\setlength{\parindent}{0pt}
\begin{tikzpicture}[x=1cm,y=1cm]

% ── Page background ─────────────────────────────────────────────────────
\begin{pgfonlayer}{page background}
  \fill[Paper,rounded corners=2mm] (-.6,-63.2) rectangle (15.6,1.2);
\end{pgfonlayer}

% ═══════════════════════════════════════════════════════════════════════
% SECTION 1: System Context
% ═══════════════════════════════════════════════════════════════════════
\SectionHeader{0.4}{PART 1 · 系统上下文}{Agent Framework Suite}

\node[font=\small,text=Slate,anchor=west,text width=14cm] at (0,-.65)
  {跨语言端到端一致性测试套件：同时驱动 Python 与 Rust 工作流引擎，验证行为相同。};

% Actors
\node[actor] (user) at (7.5, -2.0) {用户 / CI};

% Suite (central)
\node[suite] (suite) at (7.5, -4.0) {\textbf{Agent Framework Suite}\\[-1pt]{\small Go CLI · 用例目录 · 验证器}};

% Drivers
\node[driver] (py-drv) at (3.0, -6.2) {\textbf{Python Driver}\\[-1pt]{\scriptsize driver.py}};
\node[driver] (rs-drv) at (12.0, -6.2) {\textbf{Rust Driver}\\[-1pt]{\scriptsize main.rs}};

% Frameworks (external)
\node[framework] (py-fw) at (3.0, -8.4) {agent-framework\\[-1pt]{\scriptsize Python 库}};
\node[framework] (rs-fw) at (12.0, -8.4) {agent-framework-rs\\[-1pt]{\scriptsize Rust 库}};

% Catalog
\node[component] (catalog) at (12.5, -3.3) {用例目录 / Oracle\\[-1pt]{\scriptsize 10 个 WF-NNN}};

% Trust boundary
\begin{scope}[on background layer]
  \node[boundary, fit=(py-drv)(rs-drv)(py-fw)(rs-fw)] (iso) {};
\end{scope}
\node[font=\scriptsize, text=Slate!70, anchor=north west] at ($(iso.north west)+(2mm,-1mm)$) {隔离子进程边界};

% Edges
\draw[showcase/flow] (user) -- node[showcase/node,fill=white,draw=Fog,minimum height=6mm,font=\scriptsize,inner sep=1.5pt]{just check} (suite);
\draw[showcase/flow] (suite) -| node[near start,above,font=\scriptsize,text=Slate]{subprocess} (py-drv);
\draw[showcase/flow] (suite) -| node[near start,above,font=\scriptsize,text=Slate]{subprocess} (rs-drv);
\draw[showcase/softflow] (suite) -- (catalog);
\draw[showcase/softflow] (py-drv) -- node[left,font=\scriptsize,text=Slate]{公开 API} (py-fw);
\draw[showcase/softflow] (rs-drv) -- node[right,font=\scriptsize,text=Slate]{公开 API} (rs-fw);

% Parity arc
\coordinate (parity-l) at ($(py-drv.south)+(0,-.5)$);
\coordinate (parity-r) at ($(rs-drv.south)+(0,-.5)$);
\draw[showcase/data, densely dashed] (py-drv.south) -- (parity-l) -| node[below,font=\scriptsize\bfseries,text=Ocean,fill=Paper,inner sep=1pt]{parity 比对} (parity-r) -- (rs-drv.south);


% ═══════════════════════════════════════════════════════════════════════
% SECTION 2: Internal Architecture
% ═══════════════════════════════════════════════════════════════════════
\SectionHeader{-11.5}{PART 2 · 内部架构}{三层组件}

\node[font=\small,text=Slate,anchor=west,text width=14cm] at (0,-12.55)
  {自上而下：CLI 分发命令、Suite 调度用例与验证、Driver 隔离子进程调用框架。};

% CLI Layer
\begin{scope}[on background layer]
  \node[phase box, fill=Sky!6, minimum height=2.0cm, anchor=north west] at (0,-13.2) (cli-layer) {};
\end{scope}
\node[font=\scriptsize\bfseries,text=Teal,anchor=north west] at (0.3,-13.35) {CLI 层};

\node[component, fill=Sky!12, draw=Sky!60, minimum width=38mm] (main-go) at (3.5,-14.4) {\textbf{main.go}\\[-1pt]{\scriptsize flag · 命令分发 · 退出码}};
\node[font=\scriptsize,text=Slate,anchor=west] at (6.0,-14.4) {4 个子命令：doctor / list / run / version};

% Suite Layer
\begin{scope}[on background layer]
  \node[phase box, fill=Ocean!5, minimum height=3.4cm, anchor=north west] at (0,-15.6) (suite-layer) {};
\end{scope}
\node[font=\scriptsize\bfseries,text=Teal,anchor=north west] at (0.3,-15.75) {Suite 层（internal/suite）};

\node[component, fill=Ocean!10, draw=Ocean!50, minimum width=32mm] (runner) at (3.0,-16.8) {\textbf{runner.go}\\[-1pt]{\scriptsize 调度 · 验证 · parity}};
\node[component] (catalog-go) at (3.0,-18.2) {catalog.go\\[-1pt]{\scriptsize 10 个用例定义}};
\node[component] (selector-go) at (9.5,-16.8) {selector.go\\[-1pt]{\scriptsize ID / 名称 / tag}};
\node[component] (types-go) at (9.5,-18.2) {types.go\\[-1pt]{\scriptsize Observation · Report}};

% Driver Layer
\begin{scope}[on background layer]
  \node[phase box, fill=Mint!30, minimum height=2.0cm, anchor=north west] at (0,-19.4) (driver-layer) {};
\end{scope}
\node[font=\scriptsize\bfseries,text=Teal,anchor=north west] at (0.3,-19.55) {Driver 层（internal/driver）};

\node[component, fill=Mint!60, draw=Teal!40, minimum width=38mm] (client-go) at (3.5,-20.6) {\textbf{client.go}\\[-1pt]{\scriptsize 子进程 · JSON · 限流}};

\node[font=\scriptsize,text=Slate,anchor=west] at (6.0,-20.6) {临时 HOME · 固定 PATH · UTC · 64KB};

% Subprocess drivers
\coordinate (sub-sw) at (-.3,-22.6);
\coordinate (sub-ne) at (15,-21.4);
\begin{scope}[on background layer]
  \node[boundary, fit=(sub-sw)(sub-ne)] (sub-boundary) {};
\end{scope}
\node[font=\scriptsize,text=Slate!70,anchor=north west] at ($(sub-boundary.north west)+(2mm,-1.2mm)$) {隔离子进程};

\node[framework, minimum width=34mm] (py-drv2) at (3.5,-22.0) {driver.py · Python};
\node[framework, minimum width=34mm] (rs-drv2) at (11.5,-22.0) {main.rs · Rust};

% Edges between layers
\draw[showcase/flow] (main-go.south) -- (runner.north);
\draw[showcase/softflow] (runner.south) -- (catalog-go.north);
\draw[showcase/flow] (client-go.north) -- ++(0,.7) -| (runner.south east);
\draw[showcase/flow] (client-go.south west) -- ++(0,-.45) -| (py-drv2.north);
\draw[showcase/flow] (client-go.south east) -- ++(0,-.45) -| (rs-drv2.north);


% ═══════════════════════════════════════════════════════════════════════
% SECTION 3: Runtime Data Flow
% ═══════════════════════════════════════════════════════════════════════
\SectionHeader{-25.0}{PART 3 · 运行时数据流}{一个用例的生命周期}

\node[font=\small,text=Slate,anchor=west,text width=14cm] at (0,-26.05)
  {选中用例\,$\rightarrow$\,两个子进程\,$\rightarrow$\,框架调用\,$\rightarrow$\,归一化观测\,$\rightarrow$\,oracle 验证\,$\rightarrow$\,parity 比对。};

% Top row: horizontal flow
\node[suite, minimum width=26mm, minimum height=10mm, fill=Ocean!8, draw=Ocean!50] (f-select) at (1.5,-27.5) {选中用例};
\node[suite, minimum width=26mm, minimum height=10mm, fill=Sky!10, draw=Sky!50] (f-spawn) at (5.0,-27.5) {启动子进程};
\node[suite, minimum width=28mm, minimum height=10mm, fill=Fog!60, draw=Slate!40] (f-call) at (8.8,-27.5) {调用框架 API};
\node[suite, minimum width=26mm, minimum height=10mm, fill=Ocean!8, draw=Ocean!50] (f-obs) at (12.5,-27.5) {归一化观测};

% Bottom row: validation
\node[suite, minimum width=30mm, minimum height=10mm, fill=Teal!10, draw=Teal!50] (f-oracle) at (4.2,-29.8) {Oracle 验证};
\node[suite, minimum width=30mm, minimum height=10mm, fill=Ocean!10, draw=Ocean!50] (f-parity) at (8.8,-29.8) {Parity 比对};
\node[suite, minimum width=28mm, minimum height=10mm, fill=Coral!12, draw=Coral!50] (f-result) at (12.5,-29.8) {CaseResult};

% Fan-out indicator
\node[tag] at (5.0,-28.4) {×2（Py + Rs）};

% Edges
\draw[showcase/flow] (f-select) -- (f-spawn);
\draw[showcase/flow] (f-spawn) -- (f-call);
\draw[showcase/flow] (f-call) -- (f-obs);
\draw[showcase/data] (f-obs.south) -- ++(0,-.5) -| (f-oracle.north);
\draw[showcase/flow] (f-oracle) -- (f-parity);
\draw[showcase/flow] (f-parity) -- (f-result);

% Annotations
\node[font=\scriptsize,text=Slate,anchor=north,text width=3.5cm,align=center] at (f-select.south |- 0,-28.7)
  {从 catalog.go 的 10 个用例中按 selector 选中};
\node[font=\scriptsize,text=Slate,anchor=north,text width=3.5cm,align=center] at (f-call.south |- 0,-28.7)
  {driver 调用各自框架的 WorkflowBuilder 公开 API};
\node[font=\scriptsize,text=Slate,anchor=north,text width=3cm,align=center] at (f-result.south |- 0,-31.1)
  {exit 0 全部通过\\exit 1 有失败};


% ═══════════════════════════════════════════════════════════════════════
% SECTION 4: Case Matrix
% ═══════════════════════════════════════════════════════════════════════
\SectionHeader{-33.5}{PART 4 · 用例矩阵}{10 个必过用例}

\node[font=\small,text=Slate,anchor=west,text width=14cm] at (0,-34.55)
  {覆盖 workflow-core/v1 契约的全部公开能力。};

% Case boxes in two columns
\node[component, fill=Ocean!8, draw=Ocean!30, minimum width=62mm, minimum height=8mm, font=\scriptsize, anchor=west]
  (c01) at (0.2,-35.6) {\textbf{WF-001}\enspace chain\_happy\hfill core};
\node[component, fill=Ocean!8, draw=Ocean!30, minimum width=62mm, minimum height=8mm, font=\scriptsize, anchor=west]
  (c02) at (0.2,-36.6) {\textbf{WF-002}\enspace build\_missing\_start\hfill validation};
\node[component, fill=Ocean!8, draw=Ocean!30, minimum width=62mm, minimum height=8mm, font=\scriptsize, anchor=west]
  (c03) at (0.2,-37.6) {\textbf{WF-003}\enspace max\_iterations\hfill runtime};
\node[component, fill=Ocean!8, draw=Ocean!30, minimum width=62mm, minimum height=8mm, font=\scriptsize, anchor=west]
  (c04) at (0.2,-38.6) {\textbf{WF-004}\enspace fan\_out\_fan\_in\hfill routing};
\node[component, fill=Ocean!8, draw=Ocean!30, minimum width=62mm, minimum height=8mm, font=\scriptsize, anchor=west]
  (c05) at (0.2,-39.6) {\textbf{WF-005}\enspace fresh\_run\_isolation\hfill state};

\node[component, fill=Ocean!8, draw=Ocean!30, minimum width=62mm, minimum height=8mm, font=\scriptsize, anchor=west]
  (c06) at (8.0,-35.6) {\textbf{WF-006}\enspace switch\_default\_position\hfill routing};
\node[component, fill=Ocean!8, draw=Ocean!30, minimum width=62mm, minimum height=8mm, font=\scriptsize, anchor=west]
  (c07) at (8.0,-36.6) {\textbf{WF-007}\enspace checkpoint\_resume\hfill checkpoint};
\node[component, fill=Ocean!8, draw=Ocean!30, minimum width=62mm, minimum height=8mm, font=\scriptsize, anchor=west]
  (c08) at (8.0,-37.6) {\textbf{WF-008}\enspace checkpoint\_graph\_mismatch\hfill checkpoint};
\node[component, fill=Ocean!8, draw=Ocean!30, minimum width=62mm, minimum height=8mm, font=\scriptsize, anchor=west]
  (c09) at (8.0,-38.6) {\textbf{WF-009}\enspace request\_partial\_response\hfill request};
\node[component, fill=Ocean!8, draw=Ocean!30, minimum width=62mm, minimum height=8mm, font=\scriptsize, anchor=west]
  (c10) at (8.0,-39.6) {\textbf{WF-010}\enspace request\_batch\_atomicity\hfill request};

% Tag legend
\node[tag, fill=Teal!18, text=Teal, minimum width=14mm, font=\tiny\bfseries] at (2.0,-40.8) {core};
\node[tag, fill=Coral!15, text=Coral!80!black, minimum width=18mm, font=\tiny\bfseries] at (4.2,-40.8) {validation};
\node[tag, fill=Sky!18, text=Sky!70!black, minimum width=18mm, font=\tiny\bfseries] at (6.6,-40.8) {routing};
\node[tag, fill=Gold!25, text=Gold!60!black, minimum width=18mm, font=\tiny\bfseries] at (9.0,-40.8) {checkpoint};
\node[tag, fill=Indigo!15, text=Indigo, minimum width=18mm, font=\tiny\bfseries] at (11.6,-40.8) {request};


% ═══════════════════════════════════════════════════════════════════════
% SECTION 5: Verification Flow
% ═══════════════════════════════════════════════════════════════════════
\SectionHeader{-43.2}{PART 5 · 验证机制}{Oracle + Parity 双保险}

\node[font=\small,text=Slate,anchor=west,text width=14cm] at (0,-44.25)
  {每个观测先过 oracle（和预定义期望比），再做 parity（两端互比）。};

% Left: Oracle path
\node[component, fill=Teal!10, draw=Teal!50, minimum width=34mm] (obs-py) at (2.5,-45.6) {Python 观测};
\node[component, fill=Teal!10, draw=Teal!50, minimum width=34mm] (obs-rs) at (12.5,-45.6) {Rust 观测};
\node[component, fill=Teal!15, draw=Teal!60, minimum width=38mm] (oracle-box) at (7.5,-47.5) {\textbf{Oracle 验证}\\[-1pt]{\scriptsize reflect.DeepEqual}};

% Right: Parity path
\node[component, fill=Ocean!12, draw=Ocean!50, minimum width=40mm] (parity-box) at (7.5,-49.5) {\textbf{Parity 比对}\\[-1pt]{\scriptsize 去掉 participant 后逐字段}};

% Outcome
\node[component, fill=Coral!10, draw=Coral!40, minimum width=30mm] (pass-box) at (5.0,-51.3) {passed};
\node[component, fill=Coral!10, draw=Coral!40, minimum width=30mm] (fail-box) at (10.0,-51.3) {failed + proof};

% Edges
\draw[showcase/flow] (obs-py.south) -- ++(0,-.4) -| (oracle-box.north west);
\draw[showcase/flow] (obs-rs.south) -- ++(0,-.4) -| (oracle-box.north east);
\draw[showcase/flow] (oracle-box) -- node[right,font=\scriptsize,text=Slate]{两端都通过} (parity-box);
\draw[showcase/flow] (parity-box.south west) -- (pass-box.north);
\draw[showcase/error] (parity-box.south east) -- (fail-box.north);

% Normalization note
\node[font=\scriptsize,text=Slate,anchor=north west,text width=5.6cm,align=left] at (0.2,-52.3)
  {\textbf{归一化规则：}驱动抹平 UUID、时间戳、原生异常名；outcome、terminal\_state、outputs 等业务信号保留原样。};

% Isolation note
\node[font=\scriptsize,text=Slate,anchor=north west,text width=5.6cm,align=left] at (8.2,-52.3)
  {\textbf{隔离约束：}临时 HOME、固定 PATH、UTC、64KB 限流。消除环境差异。};


% ═══════════════════════════════════════════════════════════════════════
% SECTION 6: Exit Codes
% ═══════════════════════════════════════════════════════════════════════
\SectionHeader{-55.0}{PART 6 · 退出码契约}{4 种退出状态}

\node[component, fill=Teal!15, draw=Teal!50, minimum width=24mm, font=\scriptsize] (ec0) at (1.5,-56.2) {\textbf{0} 全部通过};
\node[component, fill=Coral!15, draw=Coral!50, minimum width=24mm, font=\scriptsize] (ec1) at (5.5,-56.2) {\textbf{1} 有失败};
\node[component, fill=Gold!20, draw=Gold!60!black, minimum width=24mm, font=\scriptsize] (ec2) at (9.5,-56.2) {\textbf{2} 参数错误};
\node[component, fill=Indigo!15, draw=Indigo!50, minimum width=26mm, font=\scriptsize] (ec3) at (13.5,-56.2) {\textbf{3} 驱动缺失};

% Legend
\draw[showcase/flow] (0,-58.0) -- (1.0,-58.0);
\node[legend item, anchor=west] at (1.15,-58.0) {主流程};
\draw[showcase/softflow] (3.0,-58.0) -- (4.0,-58.0);
\node[legend item, anchor=west] at (4.15,-58.0) {支撑关系};
\draw[showcase/data, densely dashed] (6.2,-58.0) -- (7.2,-58.0);
\node[legend item, anchor=west] at (7.35,-58.0) {数据流 / 比对};
\draw[showcase/error] (9.8,-58.0) -- (10.8,-58.0);
\node[legend item, anchor=west] at (10.95,-58.0) {失败路径};

\node[font=\scriptsize,text=Slate,anchor=east] at (14.8,-60.0) {workflow-core/v1 · Go 1.25 · hermetic · 无网络依赖};
\node[font=\scriptsize,text=Slate,anchor=east] at (14.8,-60.7) {生成日期：2026-08-26};

\end{tikzpicture}
\end{document}
